\section{A Learning Algorithm Subject to Fairness Constraints $\cG$}
\label{sec:learning}

In this section, we present an algorithm for training a (randomized)
classifier that satisfies false-positive subgroup fairness
simultaneously for all protected subgroups specified by a family of
group indicator functions $\cG$. All of our techniques also apply to a
statistical parity or false negative rate constraint.



Let $S$ denote a set of $n$ labeled examples
$\{z_i = (x_i, x'_i), y_i)\}_{i=1}^n$, and let $\cP$ denote the
empirical distribution over this set of examples.  Let $\cH$ be a
hypothesis class defined over both the protected and unprotected
attributes, and let $\cG$ be a collection of group indicators over the
protected attributes. We assume that $\cH$ contains a constant
classifier (which implies that there is at least one fair classifier
to be found, for any distribution).




Our goal will be to find the distribution over classifiers from $\cH$
that minimizes classification error subject to the fairness constraint
over $\cG$. We will design an iterative algorithm that, when given
access to a CSC oracle, computes an optimal randomized classifier in
polynomial time.

\iffalse Note that because minimizing classification error is NP-hard
even for simple classes $\cH$ (even absent a fairness constraint), our
goal will not be to design an algorithm that runs in polynomial time
in the worst case.  while only needing to explicitly represent a
concise state (linear in the number of timesteps the algorithm has
run, but independent of the number of hypotheses $|\cH|$ or the number
of protected groups $|\cG|$). Instead, our ultimate goal is to design
a practical algorithm. We design and analyze our algorithm imagining
that we have oracles capable of solving (unconstrained) learning and
auditing problems. These will ultimately be implemented in practice
with heuristic learning algorithms.

\noindent It will be useful to have notation for the false positive
rate of a fixed classifier $h$ on a group $g$:
\[
  \FP(h, g) = \Ex{\cP}{h(x, x') \mid y = 0 , \; g(x) = 1}
\]
and for the overall false-positive rate of $h$ (over all groups $g$):
\[
  \FP(h) = \Ex{\cP}{h(x, x')  \mid y = 0 }.
\]
% \noindent For any $\alpha\in (0, 1)$, let
% $\cG_\alpha = \{g\in \cG \mid \Pr_P[g(x) = 1, y = 0] \geq \alpha\}$.
\fi


Let $D$ denote a probability distribution over $\cH$. Consider the
following {\em Fair ERM (Empirical Risk Minimization)} problem:
\begin{align}
  &  \min_{D\in \Delta_\cH}\; \Ex{h\sim D}{err(h, \cP)}\\
  \mbox{such that } \forall g\in \cG \qquad &
\fpsize(g, \cP) \; \fpdisp(g, D, \cP) \leq \gamma.
  %                                             \Pr[y = 0 , g(x) = 1]\left(\Ex{h\sim D}{\FP(h, g)} - \Ex{h\sim D}{\FP(h)} \right) \leq \gamma \label{fair1}\\
  % & \Pr[y = 0 , g(x) = 1]\left( \Ex{h\sim D}{\FP(h)} - \Ex{h\sim D}{\FP(h, g)}\right)  \leq \gamma \label{fair2}
\end{align}
where $err(h, \cP) = \Pr_\cP[h(x, x') \neq y]$, and the quantities
$\fpsize$ and $\fpdisp$ are defined in \Cref{fp-fair}.
We will write $\OPT$ to denote the objective value at the optimum for
the Fair ERM problem, that is the minimum error achieved by a
$\gamma$-fair distribution over the class $\cH$.

% Observe that this is a linear program (with one variable for every
% classifier $h \in \cH$), and
Observe that the optimization is feasible for any distribution $\cP$:
the constant classifiers that labels all points 1 or 0 satisfy all
subgroup fairness constraints. At the moment, the number of decision
variables and constraints may be infinite (if $\cH$ and $\cG$ are infinite hypothesis classes),
but we will address this momentarily.


\begin{assumpt}[Cost-Sensitive Classification Oracle]
  We assume our algorithm has access to the cost-sensitive
  classication oracles $\CSC(\cH)$ and $\CSC(\cG)$ over the classes
  $\cH$ and $\cG$.
\end{assumpt}



Our main theoretical result is an computationally efficient
oracle-based algorithm for solving the Fair ERM problem.


\begin{restatable}{theorem}{polytime}\label{thm:polytime}
  Fix any $\nu, \delta\in (0, 1)$. Then given an input of $n$ data
  points and accuracy parameters $\nu, \delta$ and access to oracles
  $\CSC(\cH)$ and $\CSC(\cG)$, there exists an algorithm runs in
  polynomial time, and with probability at least $1 - \delta$, output
  a randomized classifier $\hat D$ such that
  $err(\hat D, \cP) \leq \OPT + \nu$, and for any $g\in \cG$, the
  fairness constraint violations satisfies
  \[
    \alpha_{FP}(g, \cP) \; \beta_{FP}(g, \hat D, \cP) \leq \gamma + O(\nu).
  \]
\end{restatable}

\iffalse %OLD VERSION
\begin{theorem}\label{thm:polytime}
  Fix any $\nu, \delta\in (0, 1)$. Then given an input of $n$ data
  points and accuracy parameters $\nu, \delta$ and access to oracles
  $\CSC(\cH)$ and $\CSC(\cG)$, there exists an algorithm runs in time
  $\poly(1/\nu, \log(1/\delta))$, and with probability at least
  $1 - \delta$, the distribution over classifiers $\hat D$ that it
  outputs satisfies $err(\hat D, \cP) \leq \OPT + 2\nu$, and for any
  $g\in \cG$, the fairness constraint violations satisfies
  \[
    \alpha_{FP}(g, \cP) \; \beta_{FP}(g, p, \cP) \leq \gamma + \frac{1 +
      2\nu}{C}.
  \]
\end{theorem}
\fi







\paragraph{Overview of our solution.} We present our solution in steps:
\begin{itemize}
\item
{\bf Step 1: Fair ERM as LP.}
First, we rewrite the Fair ERM problem as a linear program with
finitely many decision variables and constraints even when $\cH$ and $\cG$ are infinite. To do
this, we take advantage of the fact that Sauer's Lemma lets us bound
the number of labellings that any hypothesis class $\cH$ of bounded
VC dimension can induce on any fixed dataset. The LP has one variable
for each of these possible labellings, rather than one
variable for each hypothesis. Moreover, again by Sauer's Lemma,
we have one constraint for each
of the finitely many possible subgroups induced by $\cG$ on the fixed
dataset, rather than one for each of the (possibly infinitely many)
subgroups definable over arbitrary datasets. This step is important
--- it will guarantee that strong duality holds.

\item
{\bf Step 2: Formulation as Game.}
We then derive the partial Lagrangian of the LP, and note that
computing an approximately optimal solution to this LP is equivalent
to finding an approximate minmax solution for a corresponding zero-sum
game, in which the payoff function $U$ is the value of the
Lagrangian. The pure strategies of the primal or ``Learner'' player
correspond to classifiers $h \in \cH$, and the pure strategies of the
dual or ``Auditor'' player correspond to subgroups $g \in \cG$.
Intuitively, the Learner is trying to minimize the sum of the prediction
error and a fairness penalty term (given by the Lagrangian), and the
Auditor is trying to penalize the fairness violation of the Learner by
first identifying the subgroup with the greatest fairness violation and
putting all the weight on the dual variable corresponding to this
subgroup. % \sw{added; pls check}
In order to reason about convergence, we restrict the set of dual
variables to lie in a bounded set: $C$ times the probability
simplex. $C$ is a parameter that we have to set in the proof of our
theorem to give the best theoretical guarantees --- but it is also a
parameter that we will vary in the experimental section.

\iffalse
\item In order to reason about convergence, we restrict the set of
  dual variables to lie in a bounded set: a scaling of the probability
  simplex. In particular, it is useful in the analysis to view the set
  of vertices in the simplex as a finite set of pure strategies for
  the Auditor player. We show that as a function of the restriction we
  choose, the approximate minmax solution of the constrained game
  continues to give an approximately optimal solution to the fair ERM
  problem.
\fi

\item
{\bf Step 3: Best Responses as CSC.}
We observe that given a mixed strategy for the Auditor, the best response problem of the Learner corresponds to a CSC problem. Similarly, given a mixed strategy for the Learner, the best response problem of the Auditor corresponds to an auditing problem (which can be represented as a CSC problem). Hence, if we have oracles for solving CSC problems, we can compute best responses for both players, in response to arbitrary mixed strategies of their opponents.

\item
{\bf Step 4: FTPL for No-Regret.}
Finally, we show that the ability to compute best responses for each
player is sufficient to implement dynamics known to converge quickly to
equilibrium in zero-sum games.  Our algorithm has the Learner
play {\em Follow the Perturbed Leader (FTPL)\/}~\cite{KV05}, which is
a no-regret algorithm, against an Auditor who at every round best
responds to the learner's mixed strategy. By the seminal result of
\citet{FS96}, the average plays of both players converge to an
approximate equilibrium.  In order to implement this in polynomial
time, we need to represent the loss of the learner as a
low-dimensional linear optimization problem. To do so, we first define
an appropriately translated CSC problem for any mixed strategy
$\lambda$ by the Auditor, and cast it as a linear optimization problem.
\end{itemize}




\subsection{Rewriting the Fair ERM Problem}
To rewrite the Fair ERM problem, we note that even though both $\cG$
and $\cH$ can be infinite sets, the sets of possible labellings on the
data set $S$ induced by these classes are finite. More formally, we
will write $\cG(S)$ and $\cH(S)$ to denote the set of all labellings
on $S$ that are induced by $\cG$ and $\cH$ respectively, that is
\[
  \cG(S) = \{(g(x_1), \ldots , g(x_n)) \mid g\in \cG\} \qquad
  \mbox{and,} \qquad \cH(S) = \{(h(X_1), \ldots , h(X_n))\mid h\in
  \cH\}
\]

We can bound the cardinalities of $\cG(S)$ and $\cH(S)$ using
Sauer's Lemma.

\begin{lemma}[Sauer's Lemma (see e.g. \cite{KV94})]\label{lem:sauer}
  Let $S$ be a data set of size $n$.  Let $d_1 = \VCD(\cH)$ and
  $d_2 = \VCD(\cG)$ be the VC-dimensions of the two classes. Then
  \[
    |\cH(S)| \leq O\left(n^{d_1} \right) \qquad \mbox{ and } \qquad
    |\cG(S)| \leq O\left(n^{d_2} \right).
  \]
\end{lemma}


Given this observation, we can then consider an equivalent
optimization problem where the distribution $D$ is over the set of
labellings in $\cH(S)$, and the set of subgroups are defined by the
labellings in $\cG(S)$. We will view each $g$ in $\cG(S)$ as a Boolean
function.

\iffalse
  Next, we will rewrite the constraints in a way that allows us to
  reduce the problem of finding the most violated constraint to a
  cost-sensitive classification problem.

\begin{lemma}
  For each $g\in \cG_\alpha$, let
  \begin{align}
    &\Phi_+(h, g) \equiv \FP(h) \Pr[y = 0 , g(x) = 1] - \Pr[h(X) = 1 ,  y = 0 , g(x) = 1] - \gamma\\
    & \Phi_-(h, g) \equiv \Pr[h(X) = 1 ,  y = 0 , g(x) = 1] - \FP(h)\Pr[y = 0 , g(x) = 1] -\gamma
  \end{align}
  and let $\Phi_\bullet(p, g) = \Ex{h\sim D}{\Phi_\bullet(h, g)}$ for
  any $\bullet \in \{-,+\}$.
  Then for each $g\in \cG_\alpha$, a distribution $p$ over $\cH(S)$
  satisfies the constraints in \Cref{fair1,fair2} if and only if
  $\Phi_+(p, g) \leq 0$ and $\Phi_-(p,g)\leq 0$. Moreover, if
  $\Phi_+(p, g), \Phi_-(p, g) \leq \eta$, then
$    \left| \Ex{h\sim D}{\FP(h)} - \Ex{h\sim D}{\FP(h, g)} \right| \leq
    \frac{\eta}{\alpha} + \beta.$
\end{lemma}
\fi

To simplify notations, we will define the following ``fairness
violation'' functions for any $g\in \cG$ and any $h\in \cH$:
  \begin{align}
    &\Phi_+(h, g) \equiv \fpsize(g, P) \, \left(\FP(h) - \FP(h, g) \right) - \gamma\\
 % \FP(h) \Pr[y = 0 , g(x) = 1] - \Pr[h(X) = 1 ,  y = 0 , g(x) = 1] - \gamma\\
    & \Phi_-(h, g) \equiv \fpsize(g, \cP)\, \left(\FP(h, g) - \FP(h)\right) - \gamma
% \Pr[h(X) = 1 ,  y = 0 , g(x) = 1] - \FP(h)\Pr[y = 0 , g(x) = 1] -\gamma
  \end{align}
Moreover, for any distribution $D$ over $\cH$, for any sign $\bullet \in \{+,-\}$
\[
  \Phi_\bullet(D, g) = \Ex{h\sim D}{\Phi_\bullet(h, g)}.
\]
\begin{claim}\label{cl:stuff}
  For any $g\in \cG$, $h\in \cH$, and any $\nu > 0$,
\[
  \max\{\Phi_+(D, g), \Phi_-(D,g)\}\leq \nu \quad\mbox{ if and only if }\quad
  \alpha_{FP}(g, \cP)\; \beta_{FP}(g, D, \cP) \leq \gamma + \nu.
\]
\end{claim}


Thus, we
will focus on the following equivalent optimization problem.
\begin{align}
    \min_{D\in \Delta_{\cH(S)}}& \Ex{h\sim D}{err(h, \cP)}   \\
\mbox{such that for each }  g\in \cG(S):\qquad
  & \Phi_+(D, g)\leq 0 \label{f1}\\
  & \Phi_-(D,g) \leq 0 \label{f2}
\end{align}

\iffalse
\begin{align*}
  \min_{D\in \Delta_\cH}\;& \Ex{h\sim D}{err(h, \cP)}   \qquad \mbox{such that } \forall g\in \cG \quad \\
                          &  \Ex{h\sim D}{\FP(h, g)} - \Ex{h\sim D}{\FP(h)} \leq \beta\\
                          &  \Ex{h\sim D}{\FP(h)} - \Ex{h\sim D}{\FP(h, g)}  \leq \beta
\end{align*}
\fi

% \subsection{Lagrangian with Restricted Dual Space}
For each pair of constraints \eqref{f1} and \eqref{f2}, corresponding
to a group $g\in \cG(S)$, we introduce a pair of dual variables
$\lambda_g^+$ and $\lambda_g^-$. The partial Lagrangian of the linear
program is the following:
\begin{align*}
  \cL(D, \lambda) = \Ex{h\sim D}{err(h, \cP)} +  \sum_{g\in \cG(S)}
  \left( \lambda_g^+\,  \Phi_+(D, g) +  \lambda_g^- \,  \Phi_-(D,g) \right)
\end{align*}

\iffalse
\begin{claim}
  The two constant classifiers that label all points as either 0 or 1
  are feasible.
\end{claim}
\fi
By Sion's minmax theorem~\citep{sion1958}, we have
\[
  \min_{D\in \Delta_{\cH(S)}}\, \max_{\lambda\in \RR_+^{2|\cG(S)|}} \cL(p,
  \lambda) = \max_{\lambda\in \RR_+^{2|\cG(S)|}}\, \min_{D\in
    \Delta_{\cH(S)}} \cL(p, \lambda) = \OPT
\]
where $\OPT$ denotes the optimal objective value in the fair ERM
problem. Similarly, the distribution
$\arg\min_{D}\, \max_{\lambda} \cL(D, \lambda)$ corresponds to an
optimal feasible solution to the fair ERM linear program.  Thus,
finding an optimal solution for the fair ERM problem reduces to
computing a minmax solution for the Lagrangian. Our algorithms will
both compute such a minmax solution by iteratively optimizing over
both the primal variables $D$ and dual variables $\lambda$. In order
to guarantee convergence in our optimization, we will restrict the
dual space to the following bounded set:
\[
  \Lambda = \{\lambda \in \RR_+^{2|\cG(S)|} \mid \|\lambda\|_1 \leq
  C\}.
\]
where $C$ will be a parameter of our algorithm. Since $\Lambda$ is a compact
and convex set, the minmax condition continues to hold \citep{sion1958}:
\begin{equation}
  \min_{D\in \Delta_{\cH(S)}}\, \max_{\lambda\in \Lambda} \cL(D,
  \lambda) = \max_{\lambda\in \Lambda}\, \min_{D\in \Delta_{\cH(S)}}
  \cL(D, \lambda) \label{minmax}
\end{equation}

If we knew an upper bound $C$ on the $\ell_1$ norm of the optimal dual
solution, then this restriction on the dual solution would not change
the minmax solution of the program. We do not in general know such a
bound. However, we can show that even though we restrict the dual
variables to lie in a bounded set, any approximate minmax solution to
\Cref{minmax} is also an approximately optimal and approximately
feasible solution to the original fair ERM problem.
% \iffalse
% \begin{theorem}
%   Let with $C$. Then
%   \begin{equation}
%     \ROPT = \min_{D\in \Delta_{\cH(S)}}\, \max_{\lambda\in \Lambda} \cL(p,
%     \lambda) = \max_{\lambda\in \Lambda}\, \min_{D\in \Delta_{\cH(S)}}
%     \cL(p, \lambda)
%   \end{equation}
%   with $\OPT \geq \ROPT \geq \OPT - $
% \end{theorem}
% \fi


\begin{restatable}{theorem}{approxmm}\label{thm:approxmm}
  Let $(\hat D, \hat \lambda)$ be a $\nu$-approximate minmax solution
  to the $\Lambda$-bounded Lagrangian problem in the sense that
\[
  \cL(\hat D, \hat \lambda) \leq \min_{D\in \Delta_{\cH(S)}}\cL(D,
  \hat \lambda) + \nu \quad \mbox{and,} \quad \cL(\hat D, \hat
  \lambda) \geq \max_{\lambda\in \Lambda} \cL(\hat D, \lambda) - \nu.
\]
Then $err(\hat D, \cP) \leq \OPT + 2\nu$ and for any $g\in \cG(S)$,
% $\Phi_-(\hat D, g), \Phi_+(\hat D, g) \leq \frac{1 + 2\nu}{C}$, and
% therefore
\[
  \alpha_{FP}(g, \cP)\; \beta_{FP}(g, \hat D, \cP) \leq \gamma + \frac{1 + 2\nu}{C}.
\]
\iffalse % such
% that $\Pr[g(x) = 1] \geq \alpha$,
% \[
%   \left| \Ex{h\sim \hat D}{\FP(h, g)} - \Ex{h\sim\hat D}{\FP(h)}\right|
%   \leq \frac{1 + 2\nu}{C}.
% \]
\fi
\end{restatable}
% Before we prove \Cref{thm:approxmm}, we will characterize the solution
% for the maximization problem over $\lambda$.













\subsection{Zero-Sum Game Formulation}
To compute an approximate minmax solution, we will first view
\Cref{minmax} as the following two player zero-sum matrix game. The
Learner (or the minimization player) has pure strategies corresponding
to $\cH$, and the Auditor (or the maximization player) has pure
strategies corresponding to the set of vertices
$\Lambda_{\text{pure}}$ in $\Lambda$ --- more precisely, each vertex
or pure strategy either is the all zero vector or consists of a choice
of a $g \in \cG(S)$, along with the sign $+$ or $-$ that the
corresponding $g$-fairness constraint will have in the
Lagrangian. More formally, we write
\[
  \Lambda_{\text{pure}} = \{\lambda \in \Lambda \mbox{ with }
  \lambda_g^\bullet = C \mid g\in \cG(S), \bullet \in \{\pm\}\} \cup
  \{\mathbf{0}\}
\]
% where $e_g$ is the standard basis vector in coordinate $G$ in
% $|G|$-dimensional space.
Even though the number of pure strategies scales linearly with
$|\cG(S)|$, our algorithm will never need to actually represent such
vectors explicitly.  Note that any vector in $\Lambda$ can be written
as a convex combination of the maximization player's pure strategies, or
in other words: as a mixed strategy for the Auditor. For any pair of
actions $(h, \lambda)\in \cH\times \Lambda_{\text{pure}}$, the payoff
is defined as
\[
  U(h, \lambda) = err(h, \cP) + \sum_{g\in \cG(S)}
  \left(\lambda_g^+\Phi_+(h, g) + \lambda_g^-\Phi_-(h, g) \right).
\]

\begin{claim}
  Let $D\in \Delta_{\cH(S)}$ and $\lambda\in \Lambda$ such that
  $(p, \lambda)$ is a $\nu$-approximate minmax equilibrium in the
  zero-sum game defined above. Then $(p, \lambda)$ is also a
  $\nu$-approximate minmax solution for \Cref{minmax}.
\end{claim}

Our problem reduces to finding an approximate equilibrium for this
game. A key step in our solution is the ability to compute best
responses for both players in the game, which we now show can be
solved by the cost-sensitive classication (CSC) oracles.

\paragraph{Learner's best response as CSC.} Fix any mixed strategy (dual
solution) $ \lambda\in \Lambda$ of the Auditor. The Learner's
best response is given by:
\begin{equation} \label{eq:br} \argmin_{D \in \Delta_{\cH(S)}} \;
  err(h, \cP) + \sum_{g\in \cG(S)} \left( \lambda_g^+ \Phi_+(D, g) +
    \lambda_g^- \Phi_-(D, g) \right)
\end{equation}
Note that it suffices for the Learner to optimize over deterministic
classifiers $h \in \cH$, rather than distributions over
classifiers. This is because the Learner is solving a linear
optimization problem over the simplex, and so always has an optimal
solution at a vertex (i.e. a single classifier $h \in \cH$). We can
reduce this problem to one that can be solved with a single call to a
{CSC oracle}.  \iffalse A \emph{cost-sensitive classification oracle}
for the class $\cH$ does the following: given as input a set of $n$
points $\{(X_i, c_i^0, c_i^1)\}$ such that $c_i^\ell$ corresponds to
the cost for predicting label $\ell$ on point $X_i$, the algorithm
outputs an hypothesis $\hat h$ such that
\[
  \hat h \in \argmin_{h\in \cH} \sum_{i = 1}^n [h(X_i) c_i^1 + (1 -
  h(X_i)) c_i^0]
\]

We can find a best response for the Learner by making a call to the
cost-sensitive classification oracle, in which we assign \fi In
particular, we can assign costs to each example $(X_i, y_i)$ as
follows:
\begin{itemize}
% \item if $y_i = 1$, then $c_i^0 = \frac{1}{n}$ and $c_i^1 = 0$;
\item if $y_i = 1$, then $c_i^0 = 0$ and $c_i^1 = - \frac{1}{n}$;
\item otherwise, $c_i^0 = 0$ and
  \begin{align}
    c_i^1 = \frac{1}{n} &+ \frac{1}{n}\sum_{g\in \cG(S)}
                          (\lambda_g^+ - \lambda_g^-) \left(\Pr[g(x) = 1\mid y = 0] -  \mathbf{1}[g(x_i) = 1]\right)
  \end{align}
\end{itemize}

\noindent Given a fixed set of dual variables $\lambda$, we will write
$\LC(\lambda) \in \RR^n$ to denote the vector of costs for labelling
each datapoint as $1$. That is, $\LC(\lambda)$ is the vector such that
for any $i \in [n]$, $\LC(\lambda)_i = c_i^1$.

\begin{remark}
  % In our reduction, the cost of correctly classifying a point is
  % always 0, and the only thing that varies is the cost of
  % misclassifying a point.
  Note that in defining the costs above, we have translated them from
  their most natural values so that the cost of labeling any example
  with 0 is 0. In doing so, we recall that by Claim \ref{claim:csc},
  the solution to a cost-sensitive classification problem is invariant
  to translation. As we will see, this will allow us to formulate the
  learner's optimization problem as a low-dimensional linear
  optimization problem, which will be important for an efficient
  implementation of follow the perturbed leader. In particular, if we
  find a hypothesis that produces the $n$ labels
  $y = (y_1,\ldots,y_n)$ for the $n$ points in our dataset, then the
  cost of this labelling in the CSC problem is by construction
  $\langle \LC(\lambda), y \rangle$. \iffalse In this case, a
  cost-sensitive classification problem is simply a weighted agnostic
  learning problem, which is equivalent to a standard agnostic
  learning problem in which some points have been duplicated in
  proportion to their weights. Hence, our use of a cost-sensitive
  classification oracle is computationally equivalent to an agnostic
  learning oracle. We describe it directly as a cost-sensitive
  classification oracle, as this is easier to implement, as we do in
  Section \ref{sec:empirical}.  \fi
\end{remark}


\iffalse
\begin{claim}
  Given the costs specified above, the cost-sensitive classification
  oracle returns an hypothesis that solves the Learner's best response problem
  in \Cref{eq:br}.
\end{claim}
\fi



\paragraph{Auditor's best response as CSC.}
\iffalse
Next, we show the best response for the Auditor is to play a pure
strategy corresponding to the one subgroup in which the learner's
randomized classifier induces the largest disparity in the
false-positive rates. In our algorithm, we assume an \emph{auditing
  oracle} that can identify such a subgroup.
\fi


Fix any mixed strategy (primal solution) $p \in \Delta_{\cH(S)}$ of
the Learner. The Auditor's best response is given by:
\begin{equation} \label{eq:abr} \argmax_{\lambda \in \Lambda} \;
  err(D, \cP) + \sum_{g\in \cG(S)} \left( \lambda_g^+ \Phi_+(D, g) +
    \lambda_g^- \Phi_-(D, g) \right) = \argmax_{\lambda \in \Lambda}
  \sum_{g\in \cG(S)} \left( \lambda_g^+ \Phi_+(D, g) + \lambda_g^-
    \Phi_-(D, g) \right)
\end{equation}

To find the best response, consider the problem of computing
$(\hat g, \hat \bullet) = \argmax_{(g, \bullet)}\Phi_\bullet(D, g)$.
There are two cases. In the first case, $p$ is a strictly feasible primal
solution: that is $\Phi_{\hat \bullet} (D, \hat g) < 0$. In this case, the
solution to \eqref{eq:abr} sets $\lambda =
\mathbf{0}$. Otherwise, if $p$ is not strictly feasible, then by the following
\Cref{lem:dualbr} the best response is to set
$\lambda_{\hat g}^{\hat \bullet} = C$ (and all other coordinates to 0).


\begin{restatable}{lemma}{dualbr}\label{lem:dualbr}
  Fix any $\overline D\in \Delta_{\cH(S)}$ such that that
  $\max_{g\in \cG(S)} \{\Phi_+(\overline D, g), \Phi_-(\overline D,
  g)\} > 0$. Let $\lambda' \in \Lambda$ be vector with one non-zero
  coordinate $(\lambda')_{g'}^{\bullet'} = C$, where
  \[
    (g',\bullet') = \argmax_{(g, \bullet)\in \cG(S)\times \{\pm\}}
    \{\Phi_\bullet(\overline D, g) \}
  \]
  Then
  $\cL(\overline D, \lambda') \geq \max_{\lambda\in
    \Lambda}\cL(\overline D, \lambda)$.
\end{restatable}







Therefore, it suffices to solve for
$\argmax_{(g, \bullet)}\Phi_\bullet(D, g)$. We proceed by solving
$\argmax_{g} \Phi_+(D, g)$ and $\argmax_{g} \Phi_-(D, g)$ separately:
 both problems can be reduced to a cost-sensitive classification
problem. To solve for $\argmax_{g} \Phi_+(D, g)$ with a CSC oracle, we
assign costs to each example $(X_i, y_i)$ as follows:
\begin{itemize}
\item if $y_i = 1$, then $c_i^0 = 0$ and $c_i^1 = 0$;
\item otherwise, $c_i^0 = 0$ and
  \begin{align}
    c_i^1 = \frac{-1}{n}\left[ \Ex{h\sim D}{\FP(h)} - \Ex{h\sim D}{h(X_i)}  \right]
  \end{align}
\end{itemize}

To solve for $\argmax_{g} \Phi_-(D, g)$ with a CSC oracle, we assign
the same costs to each example $(X_i, y_i)$, except when $y_i = 0$,
labeling ``1'' incurs a cost of
\[
  c_i^1 = \frac{-1}{n}\left[ \Ex{h\sim D}{h(X_i)} - \Ex{h\sim
      D}{\FP(h)} \right]
\]



\subsection{Solving the Game with No-Regret Dynamics}

To compute an approximate equilibrium of the zero-sum game, we will
simulate the following \emph{no-regret dynamics} between the Learner
and the Auditor over rounds: over each of the $T$ rounds, the Learner
plays a distribution over the hypothesis class according to a
\emph{no-regret} learning algorithm (Follow the Perturbed Leader), and the Auditor plays an
approximate best response against the Learner's distribution for that
round. By the result of \cite{FS96}, the average plays of both players
over time converge to an approximate equilibrium of the game, as long
as the Learner has low regret.

\begin{theorem}[\cite{FS96}]\label{fs}
  Let $D^1, D^2 , \ldots , D^T\in \Delta_{\cH(S)}$ be a sequence of
  distributions played by the Learner, and let
  $\lambda^1, \lambda^2 , \ldots , \lambda^T \in
  \Lambda_{\textrm{pure}}$ be the Auditor's sequence of approximate
  best responses against these distributions respectively. Let
  $\overline D = \frac{1}{T} \sum_{t=1}^T D^t$ and
  $\overline \lambda = \frac{1}{T}\sum_{t=1}^T \lambda^t$ be the two
  players' empirical distributions over their strategies. Suppose that
  the regret of the Learner satisfies
  \[
    \sum_{t= 1}^T \Ex{h\sim D^t}{U(h, \lambda^t)} - \min_{h\in
      \cH(S)}\sum_{t=1}^T U(h, \lambda^t) \leq \gamma_L T \quad \mbox{
      and} \quad \max_{\lambda\in \Lambda} \sum_{t= 1}^T \Ex{h\sim
      D^t}{U(h, \lambda)} - \sum_{t=1}^T \Ex{h\sim D^t}{U(h,
      \lambda^t)} \leq \gamma_A T.
  \]
  Then $(\overline D, \overline \lambda)$ is an
  $(\gamma_L + \gamma_A)$-approximate minimax equilibrium of the game.
\end{theorem}



Our Learner will play using the Follow the Perturbed Leader (FTPL), which gives a no-regret guarantee. In order to implement FPTL, we will first need to formulate the Learner's best
response problem as a linear optimization problem over a low dimensional space. For each round $t$,
let $\overline\lambda^t = \sum_{s < t} \lambda^s$ be the vector representing the sum of the actions played by the auditor over previous rounds, and recall that
$\LC(\overline \lambda^t)$ is the cost vector given by our
cost-sensitive classification reduction. Then
the Learner's best response problem against $\overline \lambda^t$ is
the following linear optimization problem
\[
  \min_{h \in \cH(S)} \langle \LC(\overline\lambda^t), h \rangle.
\]
To run the FTPL algorithm, the Learner will optimize a ``perturbed''
version of the problem above. In particular, the Learner will play a
distribution $D^t$ over the hypothesis class $\cH(S)$ that is implicitely defined by the following sampling operation. To sample a hypothesis $h$ from $D^t$, the learner solves the following randomized optimization
problem:
\begin{equation}
  \min_{h \in \cH(S)} \langle \LC(\overline\lambda^t), h \rangle + \frac{1}{\eta}\langle
  \xi, h \rangle, \label{pbr-eq}
\end{equation}
where $\eta$ is a parameter and $\xi$ is a noise vector drawn from the
uniform distribution over $[0, 1]^n$. Note that while it is
intractable to explicitly represent the distribution $D^t$ (which has support size
scaling with $|\cH(S)|$), we can sample from $D^t$ efficiently given
access to a cost-sensitive classification oracle for $\cH$. By
instantiating the standard regret bound of FTPL for online linear
optimization (\Cref{thm:ftpl-regret}), we get the following regret
bound for the Learner.


\begin{restatable}{lemma}{learneregret}\label{learner-regret}
  Let $T$ be the time horizon for the no-regret dynamics. Let
  $D^1, \ldots, D^T$ be the sequence of distributions maintained by
  the Learner's FTPL algorithm with
  $\eta = \frac{n}{(1+C)} \sqrt{\frac{1}{\sqrt{n} T}}$, and
  $\lambda^1, \ldots, \lambda^T$ be the sequence of plays by the
  Auditor. Then
 \[
   \sum_{t= 1}^T \Ex{h\sim D^t}{U(h, \lambda^t)} - \min_{h\in
     \cH(S)}\sum_{t=1}^T U(h, \lambda^t) \leq 2 n^{1/4}
   {(1+C)} \sqrt{T}
  \]
\end{restatable}


Now we consider how the Auditor (approximately) best responds to the
distribution $D^t$. The main obstacle is that we do not have an
explicit representation for $D^t$. Thus, our first step is to
approximate $D^t$ with an explicitly represented sparse distribution
$\hat D^t$. We do that by drawing $m$ i.i.d. samples from $D^t$, and
taking the empirical distribution $\hat D^t$ over the sample. The
Auditor will best respond to this empirical distribution $\hat
D^t$. To show that any best response to $\hat D^t$ is also an
approximate best response to $D^t$, we will rely on the following
uniform convergence lemma, which bounds the difference in expected
payoff for any strategy of the auditor, when played against $D^t$ as
compared to $\hat D^t$.

\begin{restatable}{lemma}{samplingerr}\label{lem:sampling-err}
  Fix any $\xi, \delta \in (0, 1)$ and any distribution $D$ over
  $\cH(S)$.  Let $h^1, \ldots, h^m$ be $m$ i.i.d. draws from $p$, and
  $\hat D$ be the empirical distribution over the realized
  sample. Then with probability at least $1 - \delta$ over the random
  draws of $h^j$'s, the following holds,
  \[
    \max_{\lambda\in \Lambda} \left| \Ex{h\sim \hat D}{U(h, \lambda)}
      - \Ex{h\sim D}{U(h, \lambda)} \right| \leq \xi,
    % \max_{\lambda'} \Ex{h\sim D}{U(h, \lambda')} \leq \Ex{h\sim
    %   D}{U(h, \lambda)} + \xi_A + \xi \quad \mbox{ and, }\quad
    % \Ex{h\sim D}{U(h, \lambda)}
  \]
  as long as $m \geq c_0\frac{C^2\left(\ln(1/\delta) + d_2\ln(n)\right)}{\xi^2}$
  for some absolute constant $c_0$ and $d_2 = \VCD(\cG)$.
\end{restatable}


\noindent Using \Cref{lem:sampling-err}, we can derive a regret bound
for the Auditor in the no-regret dynamics.



\begin{restatable}{lemma}{auditorbr}\label{lem:auditor-br}
  Let $T$ be the time horizon for the no-regret dynamics. Let
  $D^1, \ldots, D^T$ be the sequence of distributions maintained by
  the Learner's FTPL algorithm. For each $D^t$, let $\hat D^t$ be the
  empirical distribution over $m$ i.i.d. draws from $D^t$. Let
  $\lambda^1, \ldots, \lambda^T$ be the Auditor's best responses
  against $\hat D^1, \ldots , \hat D^T$. Then with probability
  $1 - \delta$,
 \[
   \max_{\lambda\in \Lambda}\sum_{t= 1}^T \Ex{h\sim D^t}{U(h,
     \lambda)} - \sum_{t=1}^T\Ex{h \sim D^t}{ U(h, \lambda^t) }\leq T
   \sqrt{\frac{c_0 C^2 (\ln(T/\delta) + d_2 \ln(n))}{m}}
  \]
  for some absolute constant $c_0$ and $d_2 = \VCD(\cG)$.
\end{restatable}






Finally, let $\overline D$ and $\overline \lambda$ be the average of
the strategies played by the two players over the course of the
dynamics. Note that $\overline D$ is an average of many
\emph{distributions} with large support, and so $\overline D$ itself
has support size that is too large to represent explicitely. Thus, we
will again approximate $\overline D$ with a sparse distribution
$\hat D$ estimated from a sample drawn from $\overline D$. Note that
we can efficiently \emph{sample} from $\overline D$ given access to a
CSC oracle. To sample, we first uniformly randomly select a round
$t\in [T]$, and then use the CSC oracle to solve the sampling problem
defined in \eqref{pbr-eq}, with the noise random variable $\xi$
freshly sampled from its distribution. The full algorithm is described
in \Cref{alg:fairnr} and we present the proof for \Cref{thm:polytime}
below. % and we now show that with high probability, the
% output pair $(\hat D, \lambda)$ is an approximate equilibrium of the
% zero-sum game, and therefore $\hat D$ yields the desired solution for
% the Fair ERM problem.







\begin{algorithm}[h]
  \caption{\bf{FairNR: Fair No-Regret Dynamics}}
 \label{alg:fairnr}
  \begin{algorithmic}
   \STATE{\textbf{Input:} distribution $\cP$ over $n$ labelled data
      points, CSC oracles $\CSC(\cH)$ and $\CSC(\cG)$, dual bound $C$,
      and target accuracy parameter $\nu, \delta$ }

    \STATE{\textbf{Initialize}: Let
      $C = 1/\nu$,
      $\overline \lambda^0 = \mathbf{0}$,
      $\eta = \frac{n}{(1+C)} \sqrt{\frac{1}{\sqrt{n} T}}$,
      $$
      m = \frac{\left(\ln(2T/\delta)d_2\ln(n)\right) C^2 c_0
        T}{\sqrt{n} (1+C)^2 \ln(2/\delta)} \quad \mbox{and, }\quad T =
      \frac{4\sqrt{n} \ln(2/\delta)}{\nu^4}
      $$ }
    \STATE{\textbf{For} $t = 1, \ldots, T$:}

    \INDSTATE{Sample from the Learner's FTPL distribution:}
    \INDSTATE[2]{\textbf{For} $s = 1 , \ldots m$:} \INDSTATE[3]{Draw a
      random vector $\xi^s$ uniformly at random from $[0, 1]^n$}
    \INDSTATE[3]{Use the oracle $\CSC(\cH)$ to compute
      $h^{(s, t)} = \argmin_{h \in \cH(S)} \langle \LC(\overline\lambda^{(t-1)}), h
      \rangle + \frac{1}{\eta}\langle \xi^s, h \rangle$ }
    \INDSTATE[2]{Let $\hat D^t$ be the empirical distribution over
      $\{h^{s, t}\}$}
\STATE{}

\INDSTATE{Auditor best responds to $\hat D^t$:}

\INDSTATE[2]{Use the oracle $\CSC(\cG)$ to compute
  $\lambda^t = \argmax_{\lambda} \Ex{h\sim \hat D}{U(h, \lambda)}$

}

\INDSTATE{{Update}: {Let}
  $\overline \lambda^t = {\sum_{t'\leq t} \lambda^{t'}}$ }

    \STATE{Sample from the average distribution
      $\overline D = \sum_{t=1}^T D^t$: }
    \INDSTATE{\textbf{For} $s = 1 , \ldots m$:}
\INDSTATE[2]{Draw a random number $r\in [T]$ and a random vector $\xi^s$ uniformly at random from $[0, 1]^n$}
    \INDSTATE[2]{Use the oracle $\CSC(\cH)$ to compute
      $h^{(r, t)} = \argmin_{h \in \cH(S)} \langle \LC(\overline\lambda^{(r-1)}), h
      \rangle + \frac{1}{\eta}\langle \xi^s, h \rangle$ }
    \INDSTATE[1]{Let $\hat D$ be the empirical distribution over
      $\{h^{r, t}\}$}
    \STATE{\textbf{Output}: $\hat D$ as a randomized classifier}
    \end{algorithmic}
  \end{algorithm}

\begin{proof}[Proof of \Cref{thm:polytime}]
  By \Cref{thm:approxmm}, it suffices to show that with probability at
  least $1 - \delta$, $(\hat D, \overline \lambda)$ is a
  $\nu$-approximate equilibrium in the zero-sum game. As a first step,
  we will rely on \Cref{fs} to show that
  $(\overline D, \overline \lambda)$ forms an approximate equilibrium.

  By \Cref{learner-regret}, the regret of the sequence
  $D^1, \ldots , D^T$ is bounded by:
  \[
    \gamma_L = \frac{1}{T}\left[\sum_{t= 1}^T \Ex{h\sim D^t}{U(h,
        \lambda^t)} - \min_{h\in \cH(S)}\sum_{t=1}^T U(h, \lambda^t)
    \right] \leq \frac{2 n^{1/4} {(1+C)}}{\sqrt{T}}
  \]

  % Let $\gamma_A^t$ be defined as
  % \[
  %   \gamma_A^t = \max_{\lambda\in \Lambda} \left| \Ex{h\sim \hat
  %       D^t}{U(h, \lambda)} - \Ex{h\sim D^t}{U(h, \lambda)} \right|
  % \]
  % By instantiating \Cref{lem:sampling-err} and applying union bound
  % across all $T$ steps, we know with probability at least
  % $1 - \delta/2$, the following holds for all $t\in [T]$:
  % \[
  %   {\gamma_A^t} \leq \sqrt{\frac{c_0 C^2 (\ln(2T/\delta) + d_2
  %       \ln(n))}{m}}
  % \]
  % where $c_0$ is the absolute constant in \Cref{lem:sampling-err} and
  % $d_2 = \VCD(\cG)$.

  % Note that by \Cref{cor:approx-br}, the Auditor is performing a
  % $\gamma_A^t$-approximate best response at each round $t$. Then we
  % can bound the Auditor's regret as follows:
  % \begin{align*}
  %   \gamma_A = \frac{1}{T}\left[\max_{\lambda\in \Lambda} \sum_{t=
  %   1}^T \Ex{h\sim D^t}{U(h, \lambda)} - \sum_{t=1}^T \Ex{h\sim
  %   D^t}{U(h, \lambda^t)}\right] &\leq \frac{1}{T}\sum_{t=
  %                                  1}^T\left( \max_{\lambda\in \Lambda} \Ex{h\sim D^t}{U(h, \lambda)} - \Ex{h\sim
  %                                  D^t}{U(h, \lambda^t)} \right)  \\
  %                                &\leq \max_T \gamma_A^t
  % \end{align*}
  By \Cref{lem:auditor-br}, with probability $1 - \delta/2$, we have
\begin{align*}
  \gamma_A &\leq \sqrt{\frac{c_0 C^2 (\ln(2T/\delta) + d_2
        \ln(n))}{m}}
\end{align*}
  We will condition on this upper-bound event on $\gamma_A$ for the
  rest of this proof, which is the case except with probability
  $\delta/2$. By \Cref{fs}, we know that the average plays
  $(\overline D, \overline \lambda)$ form an
  $(\gamma_L+\gamma_A)$-approximate equilibrium.

  Finally, we need to bound the additional error for outputting the
  sparse approximation $\hat D$ instead of $\overline D$. We can
  directly apply \Cref{lem:sampling-err}, which implies that except
  with probability $\delta/2$, the pair $(\hat D, \overline \lambda)$
  form a $R$-approximate equilibrium, with
  \begin{align*}
    R &\leq \gamma_A + \gamma_L + \frac{\sqrt{c_0 C^2(\ln(2/\delta)
        +d_2\ln(n))}}{\sqrt{m}}
      % & \leq  \sqrt{\frac{c_0 C^2 (\ln(2T/\delta) + d_2
      %   \ln(n))}{m}}
      %   + \frac{\sqrt{c_0 C^2(\ln(2/\delta)
      %   +d_2\ln(n))}}{\sqrt{m}} + \frac{2 n^{1/4} {(1+C)}}{\sqrt{T}}\\
      % & \leq
      %   2\sqrt{\frac{c_0 C^2 (\ln(2T/\delta) + d_2
      %   \ln(n))}{m}}
      %   + \frac{2 n^{1/4} {(1+C)} \sqrt{\ln(2/\delta)}}{\sqrt{T}}
  \end{align*}
  Note that $R\leq \nu$ as long as we have $C = 1/\nu$,
  \[
    m = \frac{\left(\ln(2T/\delta)d_2\ln(n)\right) C^2 c_0 T}{\sqrt{n} (1+C)^2 \ln(2/\delta)}
 \quad
    \mbox{and, }\quad T = \frac{4\sqrt{n} \ln(2/\delta)}{\nu^4}
  \]
 This completes our proof.
\end{proof}









% \subsection{Solving the Game with Fictitious Play}
% \iffalse

% To compute an approximate minmax solution, we will first view
% \Cref{minmax} as the following two player zero-sum matrix game. The
% Learner (or the minimization player) has pure strategies corresponding
% to $\cH$, and the Auditor (or the maximization player) has pure
% strategies corresponding to the set of vertices
% $\Lambda_{\text{pure}}$ in $\Lambda$ --- more precisely, each vertex
% or pure strategy consists of a choice of a $g \in \cG_\alpha$, along
% with the coefficient $+C$ or $-C$ that the corresponding $g$-fairness
% constraint will have in the Lagrangian.  In the case of finite $\cG$,
% we can write

% \[
%   \Lambda_{\text{pure}} = \{\pm C e_{g} \mid g\in \cG_\alpha\}
% \]
% where $e_g$ is the standard basis vector in coordinate $G$ in $|G|$-dimensional space.
% (This is strictly for notational convenience --- our algorithm will never need to actually
% represent such vectors explicitly.)
% Note that any vector in $\Lambda$ can be written as a convex
% combination of maximization player's pure strategies, or in other words: as a mixed strategy for the Auditor. For any pair of actions $(h, \lambda)\in \cH\times \Lambda_{\text{pure}}$, the
% payoff is defined as
% \[
%   U(h, \lambda) = err(h, \cP) + \sum_{g\in \cG(S)}
%   \left(\lambda_g^+\Phi_+(h, g) + \lambda_g^-\Phi_-(h, g) \right).
% \]

% \begin{claim}
%   Let $D\in \Delta_{\cH(S)}$ and $\lambda\in \Lambda$ such that
%   $(D, \lambda)$ is a $\nu$-approximate minmax strategy in the
%   zero-sum game defined above. Then $(D, \lambda)$ is also a
%   $\nu$-approximate minmax solution for \Cref{minmax}.
% \end{claim}

% Our problem reduces to finding an approximate equilibrium for this
% game.  \fi
% We now give an alternative algorithm for computing an approximately optimal fair classifier. Like the algorithm given in the last section, the algorithm in this section works by simulating a game-dynamic that converges to Nash equilibrium in the zero-sum game that we derived, corresponding to the Fair ERM problem. Rather than using a no-regret dynamic, we instead use  a simple iterative procedure known as \emph{Fictitious Play}
% \citep{brown_1949}. Fictitious Play dynamics has the benefit of being more practical to implement: at each round, both players simply need to compute a single best response to the empirical play of their opponents, and this optimization requires only a single call to a CSC oracle. In contrast, the FTPL dynamic we gave in the previous section requires making many calls to a
% CSC oracle per round --- a computationally expensive process ---
% in order to find a sparse approximation to the Learner's mixed strategy at that round. The disadvantage is that Fictitious Play is only known to converge to equilibrium in the limit, rather than in a polynomial number of rounds. Nevertheless, this is the algorithm that we use in our experiments --- as we will show, it performs well on real data, despite the fact that it has weaker theoretical guarantees compared to the algorithm we presented in the last section.

%  Fictitious play proceeds in rounds, and in every round each
% player chooses a best response to his opponent's empirical history of play across previous rounds, by treating it as the mixed strategy that randomizes uniformly over the empirical history.  The seminal result of
% \cite{R51} shows that the empirical mixed strategy of the two players
% converges to Nash equilibrium in any finite, bounded zero sum game.
% \iffalse
%  To formally describe fictitious play
% in our setting, we will first characterize the best response subroutine for the
% two players.

% % \sw{the following is the details that might be deferred to the next section}
% % \ar{I think this stuff belongs here.}

% % \sw{todo: talk about how sparsity of $\lambda$ is important}
% \fi

% We can now describe the full algorithm, and state a theorem regarding its convergence.

% \begin{algorithm}[h]
%   \caption{\bf{FairFictPlay: Fair Fictitious Play}}
%  \label{alg:fairfict}
%   \begin{algorithmic}
%     \STATE{\textbf{Input:} distribution $\cP$ over the labelled data
%       points, CSC oracles $\CSC(\cH)$ and $\CSC(\cG)$ for the classes
%       $\cH(S)$ and $\cG(S)$ respectively, dual bound $C$, and number
%       of rounds $T$}

%     \STATE{\textbf{Initialize}: set $h^0$ to be some classifier in
%       $\cH$, set $\lambda^0$ to be the zero vector. Let
%       $\overline D$ and $\overline \lambda$ be the point distributions that put
%       all their mass on $h^0$ and $\lambda^0$ respectively.}

%     \STATE{\textbf{For} $t = 1, \ldots, T$:}

%     \INDSTATE{\textbf{Compute the empirical play distributions}:}

%     \INDSTATE[2]{\textbf{Let} $\overline D$ be the uniform distribution
%       over the set of classifiers $\{h^0, \ldots, h^{t-1}\}$ }

%     \INDSTATE[2]{\textbf{Let}
%       $\overline \lambda = \frac{\sum_{t'< t} \lambda^{t'}}{t}$ be the
%       auditor's empirical dual vector}


%     \INDSTATE{\textbf{Learner best responds}:
%       Use the oracle $\CSC(\cH)$ to compute
%       $h^{t} = \argmin_{h \in \cH(S)} \langle \LC(\overline\lambda), h
%       \rangle$
% \iffalse
% Call the oracle
%       $\cC$ to find a classifier
%       \[
%         h^t\in \argmin_{h\in \cH} \; err(h, \cP) + \sum_{g\in
%           \cG_\alpha} \overline \lambda_g \left(\FP(h, g) - \FP(h)
%         \right)
%       \]\fi
% }

% \INDSTATE{\textbf{Auditor best responds}: Use the oracle $\CSC(\cG)$
%   to compute
%   $\lambda^t = \argmax_{\lambda} \Ex{h\sim \overline D}{U(h,
%     \lambda)}$
% \iffalse
% Call the oracle
%     $\cA$ to find a subgroup
%     \[
%       g(\overline D) \in \argmax_{g\in \cG_\alpha} |\Ex{h\sim \overline
%         D}{\FP(h, g)} - \Ex{h\sim \overline D}{\FP(h)}|
%     \]
% \fi
% }
% \iffalse  \INDSTATE{
%     Let $\lambda^t \in \Lambda_{\text{pure}}$ with
%     $\lambda^t_{g(\overline D)} = C \; \text{sgn}\left[ \Ex{h\sim
%         \overline D}{\FP(h, g(\overline D))} - \Ex{h\sim \overline
%         D}{\FP(h)}\right]$}\fi


%     \STATE{\textbf{Output:} the final empirical distribution
%       $\overline D$ over classifiers}

%     \end{algorithmic}
%   \end{algorithm}


%   \begin{theorem}\label{thm:freakingslow}
%     Fix any constant $C$.  Suppose we run Fair Fictitious Play with
%     dual bound $C$ for $T$ rounds. Then the distribution over
%     classifiers $\overline D$ it outputs satisfies
%     $err(\overline D, \cP) \leq \OPT + 2\nu$ and for any
%     $g\in \cG_\alpha$, the fairness constraint violations satisfy
%     \[
%       \alpha_{FD}(g, \cP) \; \beta_{FD}(g, p, \cP) \leq \gamma + \frac{1 +
%         2\nu}{C},
%     \]
%     % $\Phi_-(\hat D, g), \Phi_+(\hat D, g) \leq \frac{1 + 2\nu}{C}$,
% % \[
% %   \left| \Ex{h\sim \overline D}{\FP(h, g)} - \Ex{h\sim\hat
% %       D}{\FP(h)}\right| \leq \frac{1 + 2\nu}{C},
% % \]
%   where $\nu = O\left(T^{\frac{-1}{|\cG| + |\cH|-2}}
%   \right)$ % \ar{No dependence on $C$? Why did we have to bound it then?}.
% \end{theorem}




% \begin{proof}[Proof \Cref{thm:freakingslow}]
%   By the result of \cite{R51}, the output
%   $(\overline D, \overline \lambda)$ form a $\nu$-approximate
%   equilibrium of the Lagrangian game, where
%   $$\nu = O\left(T^{\frac{-1}{|\cG| + |\cH|-2}} \right).$$
%   Our result then directly follows from \Cref{thm:approxmm}.
% \end{proof}

% \begin{remark}
%   The theoretical convergence rate in \Cref{thm:freakingslow} is quite slow, compared to the first algorithm we gave. This is because the best known convergence rate for fictitious play (Droven by
%   \citet{R51}) has an exponential dependence on the number of
%   actions of each player. However, fictitious play is known to converge much more quickly in practice. In fact, \emph{Karlin's conjecture} is that fictitious play enjoys a polynomial convergence rate\footnote{A strong version of Karlin's conjecture, in which ties can be broken adversarially, has recently been disproven \cite{DP14}. But as far as the we know, fictitious play with fixed tie-breaking rules may still converge to equilibrium at a polynomial rate.} even in the worst case. We present this algorithm because it runs much more quickly per iteration, and so is more practical to run. As we show in Section \ref{sec:empirical}, this algorithm converges quickly on real data.
% \end{remark}

%   \sw{should we argue about generalization?}  \mk{we should at least
%     discuss it and argue why it is standard}\sw{todo}

%%% Local Variables:
%%% mode: latex
%%% TeX-master: "main.tex"
%%% End:
